\documentclass[12pt,a4paper]{article}

\usepackage[a4paper,margin=1in,top=1.25in]{geometry}
\usepackage{setspace}
\usepackage{lmodern}
\usepackage{microtype}
\usepackage{amsmath,amssymb}
\usepackage{graphicx}
\usepackage{booktabs}
\usepackage{caption}
\usepackage{subcaption}
\usepackage{float}
\usepackage{xcolor}
\usepackage[colorlinks=true,linkcolor=blue,urlcolor=blue]{hyperref}

\graphicspath{{results/curves/}{baseline-v2/results/curves/}{upload/}}
\setstretch{1.08}

\title{
    \vspace{2cm}
    \textbf{\Huge Systems Engineering for Deep Learning}\\[0.5cm]
    \textbf{\Large Assignment 2}\\[0.2cm]
    \large Course Code: CS6886
}
\author{
    \textbf{\Large Aditya Girish}\\[0.2cm]
    Roll No: \texttt{EP23B048}\\[0.2cm]
    B.Tech Engineering Physics\\[0.2cm]
    IIT Madras
}
\date{September 2026}

\begin{document}

\begin{figure}[t]
    \centering
    \IfFileExists{image.png}{\includegraphics[width=0.25\linewidth]{image.png}}{}
\end{figure}

\maketitle
\thispagestyle{empty}
\newpage

\iffalse % Superseded 20-epoch baseline narrative retained only as source history.
\section{Question 1: Training Baseline}

\subsection{CIFAR-10 pipeline}

The CIFAR-10 training set contains 50,000 images and the test set contains
10,000 images, each with spatial resolution $32\times32$ and one of ten
classes. The training pipeline used the following augmentation and
normalization sequence:
\begin{enumerate}
    \item \texttt{RandomCrop(32, padding=4)},
    \item \texttt{RandomHorizontalFlip()},
    \item conversion to a tensor, and
    \item channel-wise normalization using CIFAR-10 mean
    $(0.4914, 0.4822, 0.4465)$ and standard deviation
    $(0.2470, 0.2435, 0.2616)$.
\end{enumerate}
For testing, no random augmentation was applied; images were converted to
tensors and normalized using the same mean and standard deviation. A fixed
seed of 6886 was used for reproducibility.

\subsection{Model and training strategy}

The baseline is torchvision's MobileNetV2 initialized using the public
ImageNet-1K V2 pretrained weights. The first convolutional layer was adapted
to stride 1 (rather than ImageNet's initial stride 2) to preserve useful
spatial resolution for $32\times32$ CIFAR-10 inputs. The original ImageNet
classifier was replaced by a newly initialized linear layer with 10 outputs.
The width multiplier was 1.0, dropout probability was 0.2, and the standard
torchvision BatchNorm2d settings were retained (including momentum 0.1).

The new classifier was trained for the first two epochs while the pretrained
feature extractor was frozen. The entire network was then unfrozen and
fine-tuned end-to-end for the remaining 18 epochs. Training used minibatches
of 128, SGD with Nesterov momentum (momentum 0.9), initial learning rate 0.01,
weight decay $4\times10^{-5}$, and a cosine-annealing learning-rate schedule
over 20 epochs. Cross-entropy loss was used without label smoothing.

\subsection{Results}

Table~\ref{tab:baseline-results} reports the main baseline result. The checkpoint
was selected using the fixed 5,000-image validation split; the official test set
was evaluated only after selection. The supplied checkpoint's held-out test log
records 92.39\% top-1 accuracy.

\begin{table}[H]
    \centering
    \caption{Uncompressed pretrained-MobileNetV2 baseline on CIFAR-10.}
    \label{tab:baseline-results}
    \begin{tabular}{lc}
        \toprule
        Metric & Value \\
        \midrule
        Number of training epochs & 20 \\
        Held-out test top-1 accuracy & \textbf{92.39\%} \\
        Best validation accuracy & 92.86\% \\
        Final-epoch train accuracy & 97.27\% \\
        Final-epoch validation accuracy & 92.78\% \\
        Held-out test loss & 0.2501 \\
        \bottomrule
    \end{tabular}
\end{table}

Figure~\ref{fig:baseline-curves} shows the training and test curves. The
accuracy rises quickly after the pretrained backbone is unfrozen at epoch 3:
validation accuracy increases from 63.20\% at epoch 2 to 85.88\% at epoch 3. It then
improves steadily and plateaus around 92--93\% after approximately 14 epochs.

\begin{figure}[H]
    \centering
    \includegraphics[width=0.94\linewidth]{baseline.png}
    \caption{Training and test loss and top-1 accuracy for the uncompressed
    pretrained MobileNetV2 baseline.}
    \label{fig:baseline-curves}
\end{figure}

\subsection{Failure modes and limitations}

The curves indicate mild overfitting late in training: final training accuracy
is 97.27\% while final validation accuracy is 92.78\%, a gap of approximately
4.5 percentage points. Validation loss is essentially flat late in training,
whereas training loss continues to decrease. Consequently,
substantially extending this schedule is unlikely to improve generalization
without stronger regularization or additional augmentation. A second practical
failure mode is accidental CPU execution in a hosted notebook environment;
training must be run with CUDA available and the recorded device should be
checked before a long run. Finally, this baseline uses ImageNet-pretrained
weights, so its accuracy should be interpreted as transfer-learning performance
rather than accuracy obtained by training MobileNetV2 from random
initialization.

























\fi

\section{Question 1: Training Baseline}

\subsection{(a) CIFAR-10 preparation, normalization, and augmentation}

CIFAR-10 contains 50,000 training and 10,000 test RGB images of size
$32\times32$ from ten classes. A seed-6886 permutation splits the official
training set into 45,000 training examples and a fixed 5,000-example validation
set. The validation split selects the checkpoint; the official test set is
evaluated only once after selection. Thus, test accuracy is not used for model
selection. Table~\ref{tab:q1-transforms} gives the exact transform pipeline
implemented in \texttt{src/data.py}.

\begin{table}[H]
    \centering
    \small
    \caption{CIFAR-10 preprocessing. Stochastic operations apply only to the 45,000-example training split.}
    \label{tab:q1-transforms}
    \begin{tabular}{p{0.27\linewidth}p{0.61\linewidth}}
        \toprule
        Split & Transform sequence \\
        \midrule
        Training & \texttt{RandomCrop(32, padding=4)} $\rightarrow$
        \texttt{RandomHorizontalFlip()} $\rightarrow$
        \texttt{RandAugment(num\_ops=2, magnitude=7)} $\rightarrow$
        \texttt{ToTensor()} $\rightarrow$ \texttt{Normalize} $\rightarrow$
        \texttt{RandomErasing(p=0.10)} \\
        Validation / test & \texttt{ToTensor()} $\rightarrow$ \texttt{Normalize} \\
        \bottomrule
    \end{tabular}
\end{table}

Normalization uses the CIFAR-10 channel mean $(0.4914, 0.4822, 0.4465)$ and
standard deviation $(0.2470, 0.2435, 0.2616)$. RandAugment is placed before
tensor conversion because it operates on PIL images; Random Erasing is placed
after normalization because it operates on tensors. The validation and test
paths are deterministic and share the same normalization statistics as training.

\subsection{(b) MobileNetV2 configuration and training strategy}

The baseline is torchvision MobileNetV2 with public ImageNet-1K V2 weights.
It retains MobileNetV2's width multiplier of 1.0 and inverted-residual
structure, but adapts the network for CIFAR-10 as summarized in
Table~\ref{tab:q1-model}.

\begin{table}[H]
    \centering
    \small
    \caption{MobileNetV2 configuration and changes for $32\times32$ CIFAR-10 images.}
    \label{tab:q1-model}
    \begin{tabular}{p{0.29\linewidth}p{0.25\linewidth}p{0.34\linewidth}}
        \toprule
        Component & Configuration & Reason / effect \\
        \midrule
        Initialization & ImageNet-1K V2 pretrained weights & Provides transferable low-level and mid-level features. \\
        Width multiplier & $1.0$ & Keeps the standard MobileNetV2 channel widths. \\
        Stem convolution & $3\times3$, stride changed from $2$ to $1$ & Prevents an immediate $32\rightarrow16$ reduction. Feature maps progress as $32\rightarrow32\rightarrow16\rightarrow8\rightarrow4\rightarrow2$ rather than collapsing to $1\times1$ before global pooling. \\
        Classifier & Replace ImageNet \texttt{Linear(1280,1000)} with newly initialized \texttt{Linear(1280,10)} & Matches the ten CIFAR-10 classes. \\
        Dropout & $p=0.2$ & Retained in the MobileNetV2 classifier. \\
        Batch normalization & torchvision defaults: $\epsilon=10^{-5}$, momentum $0.1$, affine parameters and running statistics enabled & Retains the pretrained BN parameterization; BN is updated during end-to-end fine-tuning. \\
        Pooling / residual blocks & Unchanged & Adaptive global average pooling converts the final $2\times2$ feature map to one value per channel. \\
        \bottomrule
    \end{tabular}
\end{table}

Only the stem stride and classifier must change to accommodate the small images.
The stride-one stem preserves spatial detail that would otherwise be discarded
at the first layer, while later native stride-two inverted-residual stages still
downsample gradually. No interpolation or input upscaling is used.

Table~\ref{tab:q1-training} lists the exact training configuration. The
classifier is warmed up while the pretrained feature extractor is frozen for
the first two epochs. At epoch 3 the entire model, including batch-normalization
parameters, is unfrozen and fine-tuned end-to-end.

\begin{table}[H]
    \centering
    \small
    \caption{Baseline optimization and regularization settings.}
    \label{tab:q1-training}
    \begin{tabular}{lc}
        \toprule
        Setting & Value \\
        \midrule
        Epochs / batch size & 60 / 128 \\
        Optimizer & SGD with Nesterov momentum \\
        Momentum / initial learning rate & $0.9$ / $0.01$ \\
        Learning-rate schedule & Cosine annealing over 60 epochs \\
        Weight decay & $1\times10^{-4}$ \\
        Loss & Cross-entropy with label smoothing $0.1$ \\
        Frozen-backbone warm-up & 2 epochs; all layers trainable from epoch 3 \\
        Reproducibility & Seed 6886; deterministic validation split \\
        \bottomrule
    \end{tabular}
\end{table}

\subsection{(c) Final accuracy, learning curves, and failure modes}

Table~\ref{tab:baseline-results} reports the result for the checkpoint selected
by best validation top-1 accuracy. Its test evaluation is on the untouched,
official 10,000-image CIFAR-10 test set. The uncompressed baseline obtains
\textbf{93.69\%} top-1 accuracy.

\begin{table}[H]
    \centering
    \small
    \caption{Uncompressed pretrained-MobileNetV2 baseline on CIFAR-10.}
    \label{tab:baseline-results}
    \begin{tabular}{lc}
        \toprule
        Metric & Value \\
        \midrule
        Checkpoint-selection criterion & Best fixed-split validation top-1 \\
        Best validation top-1 accuracy (epoch 51) & 94.38\% \\
        Final-epoch train / validation top-1 & 96.06\% / 94.34\% \\
        Held-out test loss & 0.2686 \\
        Held-out test top-1 accuracy & \textbf{93.69\%} \\
        \bottomrule
    \end{tabular}
\end{table}

\begin{figure}[H]
    \centering
    \includegraphics[width=0.94\linewidth]{baseline.png}
    \caption{Loss and top-1 accuracy curves for the 60-epoch uncompressed baseline. The validation curve is from the fixed 5,000-image split; it is not the official test curve.}
    \label{fig:baseline-curves}
\end{figure}

The main learning transition occurs when the backbone is unfrozen: validation
top-1 rises from 61.60\% at epoch 2 to 85.08\% at epoch 3. Thereafter it
improves more gradually, reaching 94.38\% at epoch 51. Late in training the
train--validation gap is about 1.72 percentage points (96.06\% versus
94.34\%), indicating mild residual overfitting despite augmentation, weight
decay, and label smoothing. Validation accuracy also fluctuates by a few tenths
of a percentage point near the optimum, so checkpoint selection rather than the
final epoch is important. Finally, the result is transfer-learning performance:
it depends on ImageNet initialization and should not be compared directly with
a MobileNetV2 trained from random initialization. In practice, accidental CPU
execution is also a reproducibility risk for this workload; the logged run used
CUDA.

\end{document}
